\chapter{Követelmény, projekt, funkcionalitás}

\section{Bevezetés}

\subsection{Cél}
\par{A dokumentum célja az elkészítendő projekt funkcióinak, korlátozásainak és magas szintű, nem szakmai definiálása, illetve a specifikációban esetlegesen felmerülő kétesen értelmezhető követelmények egyértelműsítése. Kitérünk továbbá a követelményekre és a lényeges use casekre, valamint a csapat munkájának naplószerű leírására. }

\subsection{Szakterület}
\par{A szoftver egy játék, szórakoztató célokra készül, így a szakterület a játékfejlesztés.}

\subsection{Definíciók, rövidítések}
\begin{rov}
    JDK: Java Development Kit rövidítése, a Java programozási nyelvhez szükséges fejlesztői környezet.
\end{rov}

\subsection{Hivatkozások}
\urlref{http://iit.bme.hu/}{BME IIT - Programozás alapjai 3. segédanyagok, Szoftvertechnikák segédanyagok, Szoftver projekt laboratórium feladatok}

\subsection{Összefoglalás}
\par{A dokumentum további részeiben részletesenkifejtjük a projekt követelményeit, a játék funkcióit, a felhasználók jellemzőit, valamint a projekt ütemtervét és a csapat által használt eszközöket.}

\section{Áttekintés}
\subsection{Általános áttekintés}
\par{A játék egy ablakban fut, amelyet a játékosok egy gépen keresztül használnak. Billentyűzet és egér segítségével irányítják a játékot, amelynek során különböző feladatokat kell megoldaniuk.  Az adatok tárolása helyi fájlokban történik, amelyek a játék könyvtárában helyezkednek el.}

\subsection{Funkciók}
\par{A játék egy városban játszódik, melynek utcáit egy gráf élei reprezentálják, a kereszteződések pedig az említett gráf pontjai, így kiküszöbölve Zúzmaváros területének tektonikusan aktív mivoltát.
Az utaknak számos fajtája létezik. Lehetnek többsávosak, és különböző típusúak is. Léteznek tehát sima utak, hidak és alagutak. A havazás folyamatos és egyenletes, azaz a hóvastagság az utakon és a hidakon konstans emelkedik. Zúzmaraváros szeles éghajlata miatt a felüljárók alatti pár méteres útszakasz sem menekül meg a hótakarótól, azonban az alagutakban természetesen nem hull a csapadék.}
\par{Az utak állapota a hó- és jégréteg vastagságától függ. Vékony hórétegben még zavartalanul halad a forgalom, az autók és a buszok nem csúsznak meg és nem is akadnak el.
Ugyanakkor a csapadék mellett a közlekedés hatására is dinamikusan változik az utak felülete. A hó és a jég mennyiségét belső egységekben mérjük, ahol az átszámítás alapja a sűrűség és a tömörödés mértéke. Az úttesten lévő hóréteg a rajta áthaladó forgalom hatására válik jéggé. Amennyiben egy adott útszakaszon (sávon) 5 autó áthalad, a teljes ott lévő hómennyiség jégpáncéllá áll össze. A fizikai modellezés során egy egységnyi hómennyiségből pontosan egy egységnyi jégmennyiség keletkezik. Bár a belső mennyiségi egységek megmaradnak, a hó és a jég egysége centiméterben kifejezve eltérő fizikai vastagságot takar.
Fontos kiemelni, hogy a jegesedés csak az utat használó gépjárművek (autók, buszok) hatására jön létre, a közút hókotrói pedig önmagukban nem tömörítik a havat, csak a felszerelt fej típusától függően módosítják azt.}
\par{A takarító játékosok hókotrókat irányítanak, amelyek célja az utak tisztítása és a pénzkeresés. Egy hókotró egyszerre egy sávot tud takarítani. A tisztításért járó jutalom az útszakasz súlyától függ. A hókotrók telephelyeken tudnak új eszközöket vásárolni vagy fejet cserélni. A vásárolható fejek hierarchiája és működése a következőképp alakul:}
\begin{itemize}
	\item Söprő fej: A havat és a feltört jeget a közvetlen szomszédos sávba tolja. A szomszédos sáv mindig a hókotrótól eggyel jobbra található útrészt jelenti, legyen az a külső sáv, vagy már az út széle. A söprő fej a jeget nem tördeli.
	\item Hányó fej: Ez az eszköz a havat és a feltört jeget messzire szórja. Teljesítménye igen impresszív, hiszen messzebbre tudja szórni az említett út akadályokat, mint Zúzmaraváros legtöbb sávból álló útja. 
	\item Jégtörő fej: A jégpáncélt feltöri, de nem takarítja el, így akár annak is fennáll a veszélye, hogy amennyiben a törmelék nem kerül eltávolításra egy sörpő, vagy hányófejjel (esetleg sárkánnyal) felszerelt hókotró által, akkor az autók visszatömörítik a feltört jeget tömény, veszélyes jéggé.
	\item Sószóró fej: Elolvasztja a havat és a jeget. Az olvadási idő a rétegvastagságtól függ. A sózott úton pár körig nem marad meg a hó. Működéséhez fogyóeszközként sót igényel, amit a hókotró a telephelyén tud vételezni. A szóróval felszerelt hókotró a sózás két módját különbözteti meg: sóz, nem sóz.
	\item Sárkány fej: Gázturbinával azonnal elolvasztja a havat és a jeget. Működéséhez biokerozin szükséges, amit szintén a telephelyen tud feltölteni a játékos. Természetesen ezen eszköz sem működik szükségszerűen mindig, amikor a hókotró mozog, ezzel megakadályozva a drága biokerozin elpazarlását.
\end{itemize}
\par{A játék alapvetően körökre osztott társasjáték jellegű. A játékos két szerepkör közül választhat:}
\begin{itemize}
	\item Buszvezetők: Céljuk, hogy a buszokkal a két végállomás között a lehető legtöbb kört tegyék meg.
	\item Takarítók: Irányítják a hókotró flottát, menedzselik az üzemanyagot (só, kerozin) és fejlesztik az eszközöket. Egy játékoshoz tartozhat akár több hókotró is. A hókotró flotta tagjai között van lehetőség cserélgetni a takarítófejeket, viszont csak a telephelyen. 
\end{itemize}
\par{Új hókotró vásárlásakor a játékos választhat egy alapértelmezett fejet. Ez lehet a söprő vagy jégtörő.}
\par{A játékosok Pass 'N Play módon osztoznak a játékot futtató számítógépen: amint az egyik játékos végzett, azaz például az összes hókotrójával lépett, ő már befejezte az adott kört, és következik a társ. Amennyiben ő is megtette a jónak gondolt lépéseket, az autók (NPC-k) is elvégzik lépéseiket. a lakásuk és munkahelyük között próbálnak közlekedni a legrövidebb úton. Ha az út elzáródik, a legutolsó csomópontnál várakoznak, amíg a takarítók fel nem szabadítják az utat. Szerencsére a hókotró nem tud elakadni és megcsúszni sem, így előbb utóbb minden autós eljut a munkahelyére.}
\par{A játék akkor ér véget, ha túl sok baleset történik, tehát túl sok jármű csúszott össze a jeges utakon. Ekkor az extrém időjárás következményei miatt a városban vészhelyzetet és kijárási tilalmat hirdetnek. Továbbá akkor is véget ér a játék, ha a játékban résztvevő semelyik busz sem tud lépni egy adott körben.}
\par{A játék nehézsége ízlés szerint változtatható. A játék nehézségi fokait az tükrözi, hogy azok növekedésével több autó és nagyobb baleseti ráta (összecsúszás valószínűbb) van a városban. }

\subsection{Felhasználók}
\par{A felhasználók jellemzői, tulajdonságai. Bárki játszhat a játékkal, de elsődlegesen a projekt készítői fogják használni. Két fajta felhasználó van, buszvezető és takarító szerepük szerint. }

\subsection{Korlátozások}
\par{A játék futtatásához JDK futtatókörnyezet szükséges, amelynek a telepítése a felhasználó felelőssége. A játékhoz egy mindennapi számítógép szükséges.}

\subsection{Feltételezések, kapcsolatok}
\par{A projekt nem tartalmaz ilyeneket.}


\section{Követelmények}
\subsection{Funkcionális követelmények}

\begin{longtable}{| p{2cm} | p{4cm} | p{2.7cm} | p{2cm} | p{1.5cm} | p{2.0cm}|}
    \hline
    \textbf{Azonosító} & \textbf{Leírás} & \textbf{Ellenőrzés} & \textbf{Prioritás} & \textbf{Forrás} & \textbf{Use case}\\ 
    \hline
    \endfirsthead

    \hline
    \textbf{Azonosító} & \textbf{Leírás} & \textbf{Ellenőrzés} & \textbf{Prioritás} & \textbf{Forrás} & \textbf{Use case}\\ 
    \hline
    \endhead

    
    RULE1 & %Azonosító
    Hó vastagságának növekedése& %Leírás
    Egy kör elteltével pontosan egy egységgel nő a hó vastagsága & %Ellenőrzés
    MUST & %Prioritás
    Csapat & %Forrás
    Időjárás \\ %Use-case
    \hline
    
    RULE2 & %Azonosító
    Felüljáró alatt is hullik a hó a szél miatt& %Leírás
    Egy kör elteltével a felüljáró alatti utakon egységesen nő a hó vastagsága egy egységgel & %Ellenőrzés
    MUST & %Prioritás
    Csapat & %Forrás
    Időjárás \\ %Use-case
    \hline
	
    RULE3 & %Azonosító
    Az alagutakban nem hull csapadék& %Leírás
    A föld alatti utaknak minden körben nulla a hó és jég-vastagsága & %Ellenőrzés
    MUST & %Prioritás
    Csapat & %Forrás
    Időjárás \\ %Use-case
    \hline

    RULE4 & %Azonosító
    Az úton való elakadás vizsgálata & %Leírás
    Az úton lévő hó vastagsága az elakadási határérték alatt van-e & %Ellenőrzés
    SHOULD & %Prioritás
    Csapat & %Forrás
    Jármű mozgatása \\ %Use-case
    \hline
    
    RULE5 & %Azonosító
    Az úton összecsúszhatnak a járművek & %Leírás
    Nehézségi módtól függően nagy esélye van, minél nagyobb a nehézség, annál valószínűbb. & %Ellenőrzés
    SHOULD & %Prioritás
    Csapat & %Forrás
    Időjárás \\ %Use-case
    \hline
    
    RULE6 & %Azonosító
    A hó jéggé tömörül & %Leírás
    A havas úton pontosan autó/busz haladt át, a jég vastagságának ellenőrzése, a hó eltűnése & %Ellenőrzés
    MUST & %Prioritás
    Csapat & %Forrás
    Autók vezérlése \\ %Use-case
    \hline

	
    % Második sor (példa)
    TOOLS1 & %Azonosító
    A söprő fej a havat és a feltört jeget a jobbra szomszédos útrészre tolja.& %Leírás
    A hó/ a feltört jég pontosan egy sávval kerül jobbrább, de ha jég van az nem változik & %Ellenőrzés
    SHOULD & %Prioritás
    Leírás & %Forrás
    Hókotró mozgatása \\ %Use-case
    \hline
	
    TOOLS2 & %Azonosító
    A hányó fej a havat és a feltört jeget kihányja az út legszélére.& %Leírás
    A hó/a feltört jég eltűnt az úttestről de a jeget ez a fej sem módosítja  & %Ellenőrzés
    SHOULD & %Prioritás
    Leírás & %Forrás
    Hókotró mozgatása \\ %Use-case
    \hline
	
    TOOLS3 & %Azonosító
    A jégtörő fej feltöri a jégpáncélt, de az így kialakult tört jeget nem takarítja el. & %Leírás
    A jég fel lett törve, de nem tűnt el. & %Ellenőrzés
    SHOULD & %Prioritás
    Leírás & %Forrás
    Hókotró mozgatása \\ %Use-case
    \hline

    TOOLS4 & %Azonosító
    A sószóró fej elolvasztja a havatt és a jeget, és pár körig megakadályozza annak újboli kialakulását.& %Leírás
    A hó/jég eltűnt, és a só hatásának elmúlásáig nem is jelenik meg újra. & %Ellenőrzés
    SHOULD & %Prioritás
    Leírás & %Forrás
    Hókotró mozgatása \\ %Use-case
    \hline

    TOOLS5 & %Azonosító
    A sárkány fej a gázturbinával azonnal elolvasztja a havat és a jeget, amíg van biokerozin. & %Leírás
    A hó és a jég is eltűnik az útról. Ha nincs biokerozin, akkor nem tűnik el az útról semmi. & %Ellenőrzés
    SHOULD & %Prioritás
    Leírás & %Forrás
    Hókotró mozgatása \\ %Use-case
    \hline

    PLAYERS1 & %Azonosító
    A játékos lehet buszvezető és takarító.  & %Leírás
    A felhasználó nem választhat más szerepet, de választhatja a kettő közül bármelyiket.  & %Ellenőrzés
    MUST & %Prioritás
    Csapat & %Forrás
    Járműválasztás \\ %Use-case
    \hline

    PLAYERS2 & %Azonosító
    A hókotró só/kerozintartályát csak a telephelyen lehet újratölteni. & %Leírás
    Az úton nem tud a játékos újratöltést végezni. & %Ellenőrzés
    SHOULD & %Prioritás
    Csapat & %Forrás
    Fotyóeszköz vásárlása \\ %Use-case
    \hline

    PLAYERS3 & %Azonosító
    A hókotró takarítófejeit csak a telephelyen lehet cserélni. & %Leírás
    A játékos nem tud az úton fejcserét végrehajtani. & %Ellenőrzés
    SHOULD & %Prioritás
    Csapat & %Forrás
    Fejcsere \\ %Use-case
    \hline

    PLAYERS4 & %Azonosító
    A játékos több hókotrója között lehet fejet cserélni, de csak a telephelyen. & %Leírás
    A hókotrók közötti fejcsere végrehajtása akkor és csak akkor megy végbe, ha a játékos a telephelyen van MINDKÉT hókotróval.  & %Ellenőrzés
    SHOULD & %Prioritás
    Csapat & %Forrás
    Fejcsere \\ %Use-case
    \hline

    PLAYERS5 & %Azonosító
    Új hókotró vásárlásakor az alapértelmezett fej a söprő vagy a jégtörő. & %Leírás
    A játékos pontosan csak ez a két fej közül tud választani. & %Ellenőrzés
    SHOULD & %Prioritás
    Csapat & %Forrás
    Fej vásárlása \\ %Use-case
    \hline

    RULE7 & %Azonosító
    A játék turn-based és Pass 'N Play. & %Leírás
    Ha a játékos végzett a körével, akkor már nem mozoghat többet a következő körig. A játékosok egymás után kapják meg a mozgás jogát. & %Ellenőrzés
    MUST & %Prioritás
    Csapat & %Forrás
    Jármű mozgatása \\ %Use-case
    \hline

    LOGIC8 & %Azonosító
    Miután mindkét játékos lépett, megindulnak az NPC-k, azaz az autósok. & %Leírás
    Csak akkor és csak akkor indulnak meg, amikor már mindkét játékos befejezte a körét. & %Ellenőrzés
    MUST & %Prioritás
    CSAPAT & %Forrás
    Autók vezérlése \\ %Use-case
    \hline
    
    
    LOGIC9 & %Azonosító
    Az autók a lehető legrövidebb úton próbálnak közlekedni. & %Leírás
    Nem tesznek kerülőt, akkor se, ha el fognak akadni, hiszen azt ők nem tudhatják előre. & %Ellenőrzés
    MUST & %Prioritás
    Leírás & %Forrás
    Autók vezérlése \\ %Use-case
    \hline

    LOGIC10 & %Azonosító
    Az autók elakadhatnak.  & %Leírás
    Akkor és csak akkor akadnak el, ha a hó vastagsága megüti a kritikus szintet. & %Ellenőrzés
    MUST & %Prioritás
    Leírás & %Forrás
    Autók vezérlése \\ %Use-case
    \hline

    LOGIC11 & %Azonosító
    Ha az autók előtt elzáródik az út, akkor a legutolsó csomópontnál várakoznak.  & %Leírás
    Nem fordulnak vissza, nem "repülnek át" az eltorlaszolt úton, nem tűnnek el. & %Ellenőrzés
    SHOULD & %Prioritás
    Csapat & %Forrás
    Autók vezérlése \\ %Use-case
    \hline

    RULE12 & %Azonosító
    Az elakadt autók a hókotró takarítása után kiszabadulnak. & %Leírás
    Akkor és csak akkor válnak újra mozgóképessé, ha a hókotró a megfelelő fejjel eltakarította az akadályt.  & %Ellenőrzés
    SHOULD & %Prioritás
    Leírás & %Forrás
    Autók vezérlése \\ %Use-case
    \hline

    RULE13 & %Azonosító
    A hókotró nem tud elakadni és megcsúszni sem. & %Leírás
    Nem érheti semmi probléma, ami az autókra és buszokra viszont veszélyes. & %Ellenőrzés
    SHOULD & %Prioritás
    Csapat & %Forrás
    Hókotró mozgatása \\ %Use-case
    \hline

    RULE14 & %Azonosító
    Ha túl sok baleset történik, a játék véget ér.  & %Leírás
    Ezt a fajta végzetet csak a túl sok összecsúszott jármű okozhatja, más nem.  & %Ellenőrzés
    MUST & %Prioritás
    Csapat & %Forrás
    Játék leállítása \\ %Use-case
    \hline

    RULE15 & %Azonosító
    Ha egy busz sem tud mozdulni, a játék véget ér.  & %Leírás
    Csak akkor, ha az összes busz elakadt, vagy összecsúszott egy másik busszal/ autóval.  & %Ellenőrzés
    MUST & %Prioritás
    Csapat & %Forrás
    Játék leállítása \\ %Use-case
    \hline

    RULE16 & %Azonosító
    A nehézség növelésével több autó és nagyobb baleseti ráta (összecsúszás valószínűbb) van a városban  & %Leírás
    Minél nagyobb a nehézség, annál több az autó és a baleset.  & %Ellenőrzés
    MUST & %Prioritás
    Csapat & %Forrás
    Nehézségi szint választása \\ %Use-case
    \hline
    


\end{longtable}

\subsection{Erőforrásokkal kapcsolatos követelmények}
\begin{kovetelmeny}
	{Hardver} %Azonosító
	{Alapvető} %Prioritás
	{Leírás} %Forrás
	{Fizikai} %Ellenőrzés
	{A futtatáshoz szükséges egy mindennapi, átlagos számítógép.} %Leírás
\end{kovetelmeny}
\begin{kovetelmeny}
	{Grafikus} %Azonosító
	{Alapvető} %Prioritás
	{Leírás} %Forrás
	{Fizikai} %Ellenőrzés
	{A program egy ablakban fut, tehát elengedhetetlen egy grafikus megjelenítő rendszer(monitor).} %Leírás
\end{kovetelmeny}
\begin{kovetelmeny}
	{Perifériák} %Azonosító
	{Alapvető} %Prioritás
	{Leírás} %Forrás
	{Fizikai} %Ellenőrzés
	{A játék irányításához szükségesek alapvető beviteli eszközök, mint a billentyűzet és egér.} %Leírás
\end{kovetelmeny}
\begin{kovetelmeny}
	{Java} %Azonosító
	{Alapvető} %Prioritás
	{Leírás} %Forrás
	{Parancssor} %Ellenőrzés
	{Legalább Java 11-es JDK kell a projekt futtatásához.} %Leírás
\end{kovetelmeny}

\subsection{Átadással kapcsolatos követelmények}
\begin{kovetelmeny}
    {Setup} %Azonosító
    {Alapvető} %Prioritás
    {Leírás} %Forrás
    {Parancssor} %Ellenőrzés
    {A szoftver futtatásához szükséges JDK 11 környezet letöltése a felhasználó felelőssége.} %Leírás
\end{kovetelmeny}
\begin{kovetelmeny}
    {Letöltés} %Azonosító
    {Alapvető} %Prioritás
    {Leírás} %Forrás
    {Parancssor} %Ellenőrzés
    {A projekt és maga a játék letölthető a GitHub könyvtárból, ehhez hálózati kapcsolat szükséges.} %Leírás
\end{kovetelmeny}

\subsection{Egyéb nem funkcionális követelmények}
\begin{kovetelmeny}
    {Hordozható} %Azonosító
    {Alapvető} %Prioritás
    {Leírás} %Forrás
    {Parancssor} %Ellenőrzés
    {A program minden olyan operációs rendszeren működik, ahol a Java 11 futtatása lehetséges.} %Leírás
\end{kovetelmeny}
\begin{kovetelmeny}
    {Adattárolás} %Azonosító
    {Alapvető} %Prioritás
    {Leírás} %Forrás
    {Parancssor} %Ellenőrzés
    {A program a saját könyvtárában tárolja az adatokat, nem igényel hálózati kapcsolatot a futtatáshoz.} %Leírás
\end{kovetelmeny}

\section{Lényeges use-case-ek}
\subsection{Use-case leírások}
\begin{use-case}
	%név
	{Use-case Neve}
	%rövid leírás
	{Az eset rövid leírása}
	%aktorok
	{Aktorok}
	%forgatókönyv
	Forgatókönyv \newline 
        \textbf{A.1} Alternatíva
\end{use-case}

\begin{use-case}
    %név
    {Telephelyi művelet}
    %rövid leírás
    {Absztrakt ős-eset. Olyan akciókat foglal magába, amelyeket a hókotró vezető kizárólag a telephely csomóponton hajthat végre.}
    %aktorok
    {Hókotró vezető}
    %forgatókönyv
    A játékos a körében kiválaszt egy telephelyen tartózkodó hókotrót (vagy a közös raktárat), és valamilyen menedzsment akciót kezdeményez. A rendszer ellenőrzi a jogosultságot, az egyenleget és a gépjármű pozícióját. 
\end{use-case}

\begin{use-case}
    %név
    {Hókotró vásárlása}
    %rövid leírás
    {A "Telephelyi művelet" eset leszármazottja. A játékos a telephelyen új gépjárművet vesz, és beállítja az alap fejet.}
    %aktorok
    {Hókotró vezető}
    %forgatókönyv
    A játékos kiválasztja az új hókotró vásárlása opciót. A rendszer levonja a gépjármű árát a játékos egyenlegéből. A játékos felszereli a járművet az ingyenes alapértelmezett fejjel (Söprő vagy Jégtörő). Az új hókotró a telephely csomópontján játékba lép.
\end{use-case}

\begin{use-case}
    %név
    {Fej vásárlása}
    %rövid leírás
    {A "Telephelyi művelet" eset leszármazottja. A játékos új takarítófejet vásárol a raktárba vagy a járműre.}
    %aktorok
    {Hókotró vezető}
    %forgatókönyv
    A játékos kiválaszt egy megvásárolni kívánt új fejet (Hányó, Jégtörő, Söprő, Sószóró, Sárkány). A rendszer ellenőrzi az egyenletet, levonja a vételárat, az új eszköz pedig a telephely raktárába vagy közvetlenül a kiválasztott hókotróra kerül.
\end{use-case}

\begin{use-case}
    %név
    {Fogyóeszköz vásárlása}
    %rövid leírás
    {A "Telephelyi művelet" eset leszármazottja. A játékos só és biokerozin készleteket tölt fel.}
    %aktorok
    {Hókotró vezető}
    %forgatókönyv
    A játékos megadja a megvásárolni kívánt só vagy biokerozin mennyiségét. A rendszer kiszámolja az árat, levonja az összeget a játékos egyenlegéből, és jóváírja a nyersanyagot a készletben.
\end{use-case}

\begin{use-case}
    %név
    {Fejcsere}
    %rövid leírás
    {A "Telephelyi művelet" eset leszármazottja. A játékos a raktáron lévő fejeket cserélgeti a telephelyen álló hókotróján.}
    %aktorok
    {Hókotró vezető}
    %forgatókönyv
    A játékos kiválaszt egy, a raktárban már birtokolt takarítófejet, és felszereli azt a telephelyen lévő hókotróra. A korábban a járművön lévő fej visszakerül a raktárba. A művelet ingyenes.
\end{use-case}

\begin{use-case}
    %név
    {Fogyasztó fej ki be kapcsolása}
    %rövid leírás
    {A "Jármű mozgatása" eset kiterjesztése, a hókotróra specializálva. A Sárkány és Sószóró típusú fejeknél a felhasználó ki be tudja kapcsolni ezeket a fejeket, hogy spóroljon a sóval és biokerozinnal.}
    %aktorok
    {Hókotró vezető}
    %forgatókönyv
   	Fogyóeszközös fej (Sószóró vagy Sárkány) esetén a játékos a lépés előtt eldöntheti és átállíthatja, hogy a fej be legyen-e kapcsolva, így spórolva a nyersanyaggal.
\end{use-case}

\begin{use-case}
    %név
    {Alapértelmezett fej kiválasztása}
    %rövid leírás
    {A játékos kiválasztja az induló, vagy az újonnan vásárolt hókotrójának ingyenes, alapértelmezett takarítófejét.}
    %aktorok
    {Hókotró vezető}
    %forgatókönyv
	Forgatókönyv: A rendszer felkínálja a választható ingyenes alapfejeket (Söprő fej vagy Jégtörő fej). A játékos kiválasztja az egyiket. A rendszer rögzíti a választást, és a hókotró ezzel a felszereléssel jön létre, illetve lép be a játéktérre.
\end{use-case}

\begin{use-case}
%név
{Nehézségi szint választása}
%rövid leírás
{A játékos a menüben kiválasztja a tetszőleges nehézségi szintet.}
%aktorok
{Játékos}
%forgatókönyv
A Játékos interaktál a főmenüvel és kiválasztja a kívánt nehézségi szintet ami érvényben lesz a következő játékmenet folyamán.
\end{use-case}

\begin{use-case}
%név
{Járműválasztás}
%rövid leírás
{A játékos eldöntheti, hogy hókotróval vagy busszal szeretné teljesíteni a pályát.}
%aktorok
{Játékos}
%forgatókönyv
A Játékos kiválasztja a főmenüben hogy milyen járművel szeretne lenni a következő játék során. 
\end{use-case}

\begin{use-case}
%név
{Játék indítása}
%rövid leírás
{A játékos elindítja a szimulációt a korábban kiválasztott beállításokkal. A rendszer a jármű típusától függően határozza meg a kezdőpozíciót.}
%aktorok
{Játékos}
%forgatókönyv
1. A Játékos a főmenüben a "Játék indítása" gombra kattint. \newline
2. Amennyiben a választott jármű hókotró, a rendszer a hókotró állomáson helyezi el a játékost. \newline
3. Amennyiben a választott jármű busz, a rendszer a buszgarázsban helyezi el a játékost. \newline
4. A rendszer véletlenszerű pozíciókban generálja le az Autókat a pálya úthálózatán. \newline
\end{use-case}

\begin{use-case}
%név
{Hókotrósként indítás}
%rövid leírás
{A Játék indítása során a játékos kiválasztja hogy milyen takarítófejjel szeretné kezdeni a játékot.}
%aktorok
{Hókotró vezető}
%forgatókönyv
Kiválasztja a tetszőleges fejet.
\end{use-case}

\begin{use-case}
%név
{Buszvezetőként indítás}
%rövid leírás
{A Játék indítása során a játékos kiválasztja hogy melyik útvonalon szeretne sofőrködni a játék folyamán.}
%aktorok
{Buszvezető}
%forgatókönyv
Kiválasztja a tetszőleges útvonalat.
\end{use-case}

\begin{use-case}
%név
{Jármű mozgatása}
%rövid leírás
{A játékos a kiválasztott járművet a térkép segítségével tudja mozgatni a pályán.}
%aktorok
{Játékos}
%forgatókönyv
A Játékos a térképen kijelöli az aktuális pozíciójával szomszédos csomópontok egyikét, meghatározva ezzel a jármű következő úticélját.
\end{use-case}

\begin{use-case}
    %név
    {Útvonalválasztás}
    %rövid leírás
    {A buszvezető kiválasztja a két végállomás közötti útvonalat a játék elején.}
    %aktorok
    {Buszvezető}
    %forgatókönyv
    A rendszer felkínálja a lehetséges útvonalakat. A játékos kiválasztja a számára megfelelőt, amelyen a busz közlekedni fog a játék során.
\end{use-case}

\begin{use-case}
    %név
    {Hókotró mozgatása}
    %rövid leírás
    {A "Jármű mozgatása" eset leszármazottja. A hókotró áthalad egy szomszédos csomópontba, és megtisztítja az utat.}
    %aktorok
    {Hókotró vezető}
    %forgatókönyv
	A játékos a térképen kijelöli a szomszédos csomópontot. A gépjármű átmozog rajta, miközben a felszerelt fej automatikusan manipulálja a havat és a jeget. A rendszer kiszámítja az útszakasz súlya alapján járó jutalmat.
\end{use-case}

\begin{use-case}
    %név
    {Busz mozgatása}
    %rövid leírás
    {A "Jármű mozgatása" eset leszármazottja. A busz a kijelölt útvonalán halad előre egy szomszédos csomópontba.}
    %aktorok
    {Buszvezető}
    %forgatókönyv
    A játékos a térképen kijelöli a szomszédos csomópontot a menetrendje szerint. A busz átmozog rajta, közeledve a célállomás felé.
\end{use-case}

\begin{use-case}
%név
{Időjárás vezérlése}
%rövid leírás
{A játék során az időjárással kapcsolatos dolgok.}
%aktorok
{Rendszer}
%forgatókönyv
A hó egyenletesen hullik a pályán. A hó tömörödése következtében jéggé válik.
\end{use-case}

\begin{use-case}
%név
{Autók vezérlése}
%rövid leírás
{Az NPC autók irányítását és generálását végzi.}
%aktorok
{Rendszer}
%forgatókönyv
Létrejönnek autóka majd azok a kezdőpontjukból a céljukba haladnak.
\end{use-case}

\begin{use-case}
%név
{Játék leállítása}
%rövid leírás
{Figyeli a balesetek számát.}
%aktorok
{Rendszer}
%forgatókönyv
Leállítja a játékot ha túl sok baleset történik.
\end{use-case}

\subsection{Use-case diagram}
\diagram{docs/img/use_case_uml}{Use-case diagram}{16cm}

\section{Szótár}
%szotar.sh script, lsd. a dokumentációt!!!
\begin{szotar}
	\szotaritem{Alagút }{  Olyan speciális úttípus, amelyben a hó és a jég nem hullik a csapadékmentesség miatt.
 }
\szotaritem{Autó (NPC) }{  A város lakóit szállító, önműködő jármű, amely a legrövidebb úton közlekedik a munkahely és a lakás között.
 }
\szotaritem{Baleseti ráta }{  A nehézségi foktól függő mutatószám, amely meghatározza a gépjárművek jeges úton történő összeütközésének valószínűségét.
 }
\szotaritem{Biokerozin }{  A sárkány fej működtetéséhez szükséges speciális üzemanyag, amely a telephelyen tölthető fel.
 }
\szotaritem{Buszvezető }{  A játékos egyik választható szerepköre, akinek célja a legtöbb kör megtétele a végállomások között.
 }
\szotaritem{Csomópont }{  A városi úthálózatot reprezentáló gráf azon pontjai, ahol az utak (élek) találkoznak, és ahol a forgalom várakozhat.
 }
\szotaritem{Egységnyi hó/jég }{  A csapadék belső mennyiségi egysége, amely a sűrűségtől függően eltérő fizikai vastagságot (centimétert) jelenthet az úton.
 }
\szotaritem{Feltört jég }{  A jégtörő fej által szétzúzott, de el nem takarított jégtörmelék, amely az autók áthaladása miatt újra veszélyes jégpáncéllá állhat össze.
 }
\szotaritem{Hányó fej }{  Olyan nagyteljesítményű takarítóeszköz, amely a havat és a feltört jeget az úttól messzire repíti.
 }
\szotaritem{Híd }{  Olyan úttípus, amelyen az alagúttal ellentétben folyamatosan és egyenletesen gyarapodik a hóréteg.
 }
\szotaritem{Hó }{  Az utakon folyamatosan és egyenletesen gyűlő csapadék, amely a gépjárművek súlya alatt jéggé tömörödik.
 }
\szotaritem{Hókotró }{  A takarító játékos által irányított jármű, amely különböző fejekkel felszerelve tisztítja az utakat.
 }
\szotaritem{Jégpáncél }{  Az autók és buszok áthaladása miatt tömörödött hóréteg, amely csúszásveszélyt okoz.
 }
\szotaritem{Jégtörő fej }{  A hókotróra szerelhető eszköz, amely a szilárd jégpáncélt darabokra töri, de nem távolítja el az útról.
 }
\szotaritem{Kijárási tilalom }{  A játék egyik végállapota (vereség), amely akkor következik be, ha a jeges utakon túl sok baleset történik a városban.
 }
\szotaritem{Körökre osztott működés }{  A játékmenet alapja, ahol a játékosok és az NPC-k egymás után, meghatározott sorrendben teszik meg lépéseiket.
 }
\szotaritem{Olvadási idő }{  A sózott útszakaszokon a hó és jég eltűnéséhez szükséges időtartam, amely egyenesen arányos a rétegvastagsággal.
 }
\szotaritem{Pass 'N Play }{  Olyan lokális többjátékos mód, amelyben a játékosok ugyanazon a számítógépen, egymást váltva hajtják végre a köreiket.
 }
\szotaritem{Sárkány fej }{  Gázturbinával felszerelt eszköz, amely biokerozin égetésével azonnal elolvasztja a havat és a jeget.
 }
\szotaritem{Sáv }{  Az utak belső felosztása; a hókotró egyszerre csak egy ilyen egységet tud megtisztítani, a forgalom pedig ezeken halad keresztül.
 }
\szotaritem{Só }{  A sószóró fej működéséhez szükséges fogyóeszköz, amely megakadályozza a hó megmaradását az úton.
 }
\szotaritem{Söprő fej }{  Olyan alapvető takarítóeszköz, amely a havat és a feltört jeget a menetirány szerinti jobbra szomszédos sávba tolja.
 }
\szotaritem{Sószóró fej }{  Olyan eszköz, amely só alkalmazásával olvasztja el a csapadékot és ideiglenesen tisztán tartja az útfelületet.
 }
\szotaritem{Súly (útszakaszé) }{  Az egyes útszakaszokhoz rendelt érték, amely meghatározza a takarításért kapható jutalom mértékét.
 }
\szotaritem{Takarító }{  A játékos azon szerepköre, amelyben a hókotró flottát és az üzemanyag-készleteket menedzseli.
 }
\szotaritem{Telephely }{  A városban található bázis, ahol a hókotrók felszereléseit cserélni, újakat vásárolni és üzemanyagot vételezni lehet.
 }
\szotaritem{Úthálózat }{  A várost reprezentáló gráf, ahol a csomópontok a kereszteződéseket, az élek pedig a különböző típusú utakat jelentik.
 }
\szotaritem{Végállomás }{  A buszok útvonalának két végpontja, amelyek között a buszvezetőknek a lehető legtöbb kört kell teljesíteniük. }

\end{szotar}

\section{Projekt terv}
\par{A projekt elkészítése során a csapat a következő ütemtervet követi, amelyben a fontosabb mérföldkövek és határidők szerepelnek.}

\subsection{Projektütemterv}
\begin{terv}
		\tervitem{márc. 2.}  {Követelmény, projekt, funkcionalitás}{ Csapat }
        \tervitem{márc. 9. } { Analízis modell (I. változat) }{ Csapat } 
        \tervitem{márc. 16. } { Analízis modell (II. változat) }{ Csapat }
        \tervitem{márc. 23. }{ Szkeleton tervezése }{ Csapat }
        \tervitem{márc. 30. }{ Szkeleton elkészítése }{ Csapat }
        \tervitem{ápr. 13. }{ Prototípus koncepciója }{ Csapat }
        \tervitem{ápr. 20. }{ Részletes tervek }{ Csapat }
        \tervitem{máj. 4. }{ Prototípus elkészítése }{ Csapat }
        \tervitem{máj. 11. }{ Grafikus változat tervei }{ Csapat }
        \tervitem{máj. 27. }{ Grafikus változat elkészítése }{ Csapat } 
		\tervitem{máj. 29. }{ Egyesített dokumentáció }{ Csapat }
\end{terv}

\subsection{Erőforrások, eszközök}
A fejlesztés során felhasznált segédeszközök:
\begin{itemize}
	\item Dokumentáció: LaTeX
	\item Kommunikáció: Discord, Messenger
	\item Modellező eszköz: PlantUML
	\item Fejlesztő környezetek: Visual Studio Code
	\item Forráskód megosztás, verziókezelés: GitHub, Git
\end{itemize}
